% !TEX program = xelatex
\documentclass[UTF8,a4paper,11pt]{ctexart}

\usepackage[a4paper,margin=2.05cm,headheight=14pt]{geometry}
\usepackage{amsmath,amssymb,mathtools,bm}
\usepackage{booktabs,longtable,array,tabularx,multirow}
\usepackage{graphicx,xcolor,listings,xurl,hyperref,pgfplots,seqsplit}
\usepgfplotslibrary{groupplots}
\pgfplotsset{compat=1.18}

\definecolor{pyblue}{RGB}{36,91,163}
\definecolor{qudared}{RGB}{190,65,55}
\definecolor{stagegray}{RGB}{90,105,115}
\hypersetup{
  colorlinks=true,linkcolor=pyblue,urlcolor=pyblue,citecolor=pyblue,
  pdftitle={PyQCU 分布式 Strict-MG 最终 debug、优化与全矩阵对照}}
\urlstyle{tt}
\Urlmuskip=0mu plus 3mu
\newcommand{\pathref}[1]{\begingroup\Urlmuskip=0mu plus 3mu\path{#1}\endgroup}
\newcommand{\code}[1]{%
  \begingroup\ttfamily
  \expandafter\seqsplit\expandafter{\detokenize{#1}}%
  \endgroup}
\newcolumntype{L}[1]{>{\raggedright\arraybackslash}p{#1}}
\newcolumntype{Y}{>{\raggedright\arraybackslash}X}
\newcommand{\norm}[1]{\left\lVert#1\right\rVert}
\newcommand{\Yhat}{\widehat Y}
\setlength{\parindent}{2em}
\setlength{\parskip}{0.35em}
\emergencystretch=3em

\title{PyQCU 分布式 Strict MultiGrid：\\
双 P100 illegal-memory debug、并行优化与全矩阵 QUDA 对照}
\author{PyQCU 工程分析记录}
\date{2026 年 9 月 21 日}

\begin{document}
\maketitle

\begin{abstract}
本文记录 PyQCU 分布式 Strict-MG 从双 P100 旧版 clover-gauge MPI halo
非法访问修复，到真正 rank-sliced HDF5 读取、粗层通信优化，再到
\(2\times2\times3\times2\times3=72\) 个公平对照单元
（共 \(144\) 个 side case）的完整测试。最终矩阵覆盖 c64/c128、
\(8^3\times16\)、\(16^4\)、\(16\times32\times32\times48\)（c64）与
\(16\times16\times32\times32\)（c128）、cold \(+\) 2 warmup \(+\) 5 steady、
trace on/off、BiCGstab/MG-2/MG-3、V100 单卡与双 P100。
所有 \(144\) 个 side case 均通过真残差门禁，\(72\) 个比较均
\code{fair=true}。MG-2/3 的 \(48\) 个比较中，PyQCU 在 \(36\) 个单元更快，
QUDA/PyQCU 比值中位数 \(1.298\)；V100 的 MG-2/3 中位数 \(2.492\)，
双 P100 的中位数 \(1.063\)。BiCGstab reference 的 \(24\) 个比较
全部慢于 QUDA，中位数 \(0.243\)，因此报告中严格分列，不用 MG 结果
掩盖 reference 路径的不足。本文同时记录两个失败的粗层归约实验
（triple-dot 与 pipelined BiCGStab）及其回退，避免把数值失稳路径
写成性能优化。
\end{abstract}

\tableofcontents

\section{范围、验收与证据等级}

\begin{table}[htbp]
\centering
\small
\caption{最终验收结果}
\begin{tabularx}{\textwidth}{@{}L{0.32\textwidth}Y@{}}
\toprule
项目 & 实测结果 \\
\midrule
笛卡尔 side case & \(144/144\) 个单元存在且两侧 \code{status=ok} \\
公平比较 & \(72/72\) 个 \code{fair=true}，无 config hash、input bundle 或残差门禁失败 \\
相位与层级记录 & 432 个 phase cell、336 个 MG level cell 全覆盖 \\
真残差 & PyQCU 最大 \(8.01\times10^{-7}\)；QUDA 最大 \(7.97\times10^{-7}\)，全部低于各自 gate \\
分布式输入 & 双 rank 下 gauge/source/null 均采用 rank-sliced HDF5 hyperslab；round-trip 逐元素误差 0 \\
旧版 CUDA illegal access & 原失败路径不再复现；随后完成全部双 P100 长测 \\
性能结论 & MG-2/3 \(36/48\) 胜；V100 中位数 \(2.492\times\)，双 P100 中位数 \(1.063\times\) \\
BiCGstab reference & \(0/24\) 胜，中位数 \(0.243\times\)，单独列为短板 \\
\bottomrule
\end{tabularx}
\end{table}

证据等级如下：
\begin{itemize}
\item E1：本机 CUDA/MPI 实测、JSON/trace/validate 输出与 CUDA sanitizer；
\item E2：QUDA 1.1.0 与 PyQCU 源码逐层对照、算法推导；
\item E3：仅用于历史背景的旧报告结论。
\end{itemize}

最终机器可读产物为
\pathref{data/mg_matrix_final_combined/units/}、
\pathref{data/mg_matrix_final_analysis_full.json}、
\pathref{data/mg_matrix_final_validation_full.json}、
\pathref{data/mg_matrix_final_report/mg_matrix.csv} 与
\pathref{data/mg_matrix_final_report/mg_levels.csv}。

\section{双 P100 非法访问的根因与修复}

\subsection{设备绑定根因}

旧版双 P100 收集器在 \code{CUDA_VISIBLE_DEVICES=0,1} 下由 Python
选择 Torch 逻辑设备 0/1，而 C++ \code{LatticeSet} 仍强制
\code{cudaSetDevice(0)}。rank 1 因此把 GPU1 张量指针交给 GPU0
上下文，在旧版 clover-gauge \code{pick_up_u_*} pack kernel 中触发
CUDA illegal memory access。修复分两处：
\begin{itemize}
\item \pathref{cpp/cuda/qcu/include/define.h:271} 将
  \code{_TEST_SINGLE_GPU_MULTI_RANK_} 设为 0；
\item \pathref{cpp/cuda/qcu/include/lattice_set.h:272--288} 在多 rank
  路径读取 \code{PYQCU_MPI_DEVICE_ID}，回退 local rank；
  \pathref{examples/qcu/dev87/bench_strict_vs_quda.py:3181} 把 Torch
  实际设备索引写入该环境变量。
\end{itemize}

\subsection{pack kernel 边界根因}

旧版 10 个 \code{pick_up_u_*}（x/y/z/t、xy/xz/xt/yz/yt/zt）pack kernel
在线程数按 face grid 向上取整时没有上界检查；尾线程越过 face 元素数。
\pathref{cpp/cuda/qcu/src/clover_dslash_comm.cu} 与
\pathref{cpp/cuda/qcu/include/lattice_clover_dslash.h} 中的所有
1D/2D pack kernel 现在在索引解码前返回越界线程。修复后的原失败命令
退出码为 0；compute-sanitizer memcheck 未再报告 illegal access。

\section{真正的分布式输入与粗层通信优化}

\subsection{rank-sliced HDF5}

旧 worker 每个 MPI rank 读取完整 gauge/source/null HDF5，因此双 rank
重复占用数 GiB host memory。新实现
\pathref{examples/qcu/dev87/bench_strict_vs_quda.py:2520} 的
\code{_load_h5_local_array} 使用 h5py MPI-IO hyperslab：
\begin{itemize}
\item gauge/source 的 checkerboard-compressed 数据集按
  \(x,y,z,t/2\) 读取 rank slab；局部原点和长度均为偶数，所以局部
  checkerboard 与全局 checkerboard 一致；
\item canonical full null-vector 按 \(x,y,z,t\) 读取 rank slab；
\item 单 rank 保留全量读取语义；多 rank 对 dtype/shape/file-size
  重新校验，父进程仍在启动前计算完整 logical SHA-256。
\end{itemize}
双 rank 小格实测：gauge、source、null 三类 slab 与全量读取切片
逐元素 \code{max_abs=0}。

\subsection{Strict-MG 数据路径优化}

\begin{enumerate}
\item \textbf{活跃 halo slot staging.}
  \pathref{cpp/cuda/qcu/src/apply_multigrid_strict.cu:1690--1760}
  仅复制 process grid 中实际活跃的 \((\mathrm{dim},\mathrm{side})\)
  slot，而不是无条件复制完整 8-slot 缓冲；
\item \textbf{link halo cache.}
  \pathref{cpp/cuda/qcu/src/apply_multigrid_strict.cu:1808--1835}
  按 \((parity,dim)\) 缓存 static preconditioned link ghost，
  pointer 不变时跳过重复 pack/MPI；
\item \textbf{prolongation 与加法融合.}
  \pathref{cpp/cuda/qcu/src/apply_multigrid_strict.cu:1489--1531}
  用 \code{add} 参数在同一 kernel 中完成 prolongation 加 corrected
  coarse contribution；
\item \textbf{trace-on cache 复用.}
  \pathref{examples/qcu/dev87/bench_strict_vs_quda.py:6135} 允许显式
  \code{cache-expect=hit} 时 cold/main 两个 phase 复用同一个
  hierarchy cache；
\item \textbf{矩阵编排.}
  \pathref{examples/qcu/dev87/bench_mg_matrix_full.py:394} 记录
  \code{strict_cache_expect}，trace-on 优先复用对应 trace-off cache。
\end{enumerate}

\section{QUDA 的优势与 PyQCU 的剩余短板}

\subsection{QUDA 的结构性优势}

\begin{itemize}
\item QUDA 的 transfer、block orthogonalization、prolongator/restrictor、
  coarse \(X/Y/\Yhat\) 与 coarse-PC 均有成熟专用 kernel；
\item QUDA 的粗层求解器默认使用 CA-GCR，一次 block dot 可摊销多个
  Krylov 方向的归约延迟；
\item QUDA 有完整的 QMP/NVSHMEM/coarse halo、逐层精度和长期硬件回归；
\item QUDA 的 coarse dslash 与通信路径可直接复用生产 kernel，而
  PyQCU Strict 的粗层仍是“宿主侧全局归约 + halo 交换”的较薄实现。
\end{itemize}

\subsection{PyQCU 的剩余短板}

\begin{itemize}
\item 最粗层的每次 BiCGStab 迭代仍需 host \(\to\) device \(\to\)
  host 标量回传；双 P100 上粗层延迟是首要瓶颈；
\item 旧版 Strict 的 mixed precision、奇数 coarse extent 和
  non-nearest-neighbour stencil 仍有 fail-closed 限制；
\item Galerkin setup 与 runtime cache 虽已分离，但 cache miss 仍是
  分钟级成本；矩阵编排已通过 trace-off/trace-on 共享 cache 降低重复；
\item BiCGstab reference 没有 MG 的多层预处理，单次迭代成本高，
  在 24 个 L1 对照中全部慢于 QUDA。
\end{itemize}

\subsection{两条失败优化路径}

为降低粗层归约延迟，曾实现并测试：
\begin{enumerate}
\item 将 \(\langle s,s\rangle,\langle t,s\rangle,\langle t,t\rangle\)
  合并成一个 triple reduction；
\item 将 BiCGStab 的 \(\rho\) 延迟到残差更新后，与
  \(\langle r,r\rangle\) 合成 dot-pair，形成 pipelined 版本。
\end{enumerate}
两者在小格、中格能收敛，但在
\(16\times32\times32\times48\)、c64、双 P100 的 FGMRES 外循环中
出现残差停滞；trace 显示外残差长期停在 \(10^{-3}\) 量级而 coarse
residual 仍约 \(0.2\)。这属于当前分布式 BiCGStab 对归约重排的数值
敏感性，不是可接受的性能优化。两路径均已从最终源码移除，正式
结果使用原始、已验证的归约顺序。该失败实验说明“减少同步次数”不能
以牺牲粗层递推稳定性为代价。

\section{测试协议与覆盖}

\begin{table}[htbp]
\centering
\small
\caption{最终矩阵协议}
\begin{tabularx}{\textwidth}{@{}L{0.25\textwidth}Y@{}}
\toprule
轴 & 取值 \\
\midrule
side & PyQCU、QUDA \\
precision & c64、c128  \\
c64 lattice & \(8^3\times16\)、\(16^4\)、\(16\times32\times32\times48\) \\
c128 lattice & \(8^3\times16\)、\(16^4\)、\(16\times16\times32\times32\) \\
trace & off、on \\
levels & BiCGstab reference、MG-2、MG-3（执行顺序 2,3,1） \\
phases & cold \(1\)、warmup \(2\)、steady \(5\) \\
solver gate & c64 \(5\times10^{-6}\)、c128 \(5\times10^{-8}\) true residual \\
device & V100-SXM2-32GB 单卡；P100-PCIE-16GB \(\times2\)，process grid \([2,1,1,1]\) \\
shared parameters & coarse-tol \(0.3\)；P100 全组 nu=1；V100 c64 nu=1；V100 c128 nu=2 \\
\bottomrule
\end{tabularx}
\end{table}

矩阵共 \(144\) 个 side case。每个 PyQCU/QUDA 对共享同一 canonical
null-vector、同一 gauge、同一 solver tolerance、同一 restart 与同一
coarse-tol。两个侧别均记录 phase、iterations、true residual，
trace-on 侧额外记录逐层 stage。

\section{性能结果}

\subsection{总览}

\begin{table}[htbp]
\centering
\small
\caption{MG-2/3 分组结果；比值为 QUDA/PyQCU，\(>1\) 表示 PyQCU 更快。
每组 6 个单元（3 lattice \(\times\) 2 trace）。}
\begin{tabular}{llrrrrr}
\toprule
device & precision & levels & median & wins & losses & range \\
\midrule
P100 & c64 & MG-2 & 0.923 & 1 & 5 & 0.645--1.110 \\
P100 & c64 & MG-3 & 1.157 & 4 & 2 & 0.640--1.777 \\
P100 & c128 & MG-2 & 1.072 & 4 & 2 & 0.823--1.294 \\
P100 & c128 & MG-3 & 1.098 & 4 & 2 & 0.661--1.365 \\
V100 & c64 & MG-2 & 2.212 & 5 & 1 & 0.999--4.446 \\
V100 & c64 & MG-3 & 2.184 & 6 & 0 & 1.207--9.115 \\
V100 & c128 & MG-2 & 3.012 & 6 & 0 & 1.405--16.393 \\
V100 & c128 & MG-3 & 3.795 & 6 & 0 & 1.602--19.587 \\
\bottomrule
\end{tabular}
\end{table}

\begin{table}[htbp]
\centering
\small
\caption{汇总统计}
\begin{tabular}{lrrrrr}
\toprule
集合 & cases & median & wins & losses & range \\
\midrule
MG-2/3 全体 & 48 & 1.298 & 36 & 12 & 0.640--19.587 \\
BiCGstab reference & 24 & 0.243 & 0 & 24 & 0.164--0.468 \\
V100 MG-2/3 & 24 & 2.492 & 23 & 1 & 1.000--19.587 \\
P100 MG-2/3 & 24 & 1.063 & 13 & 11 & 0.640--1.777 \\
\bottomrule
\end{tabular}
\end{table}

代表性单元：
\begin{itemize}
\item 双 P100、\(16\times32\times32\times48\)、c64、MG-2、trace-off：
  PyQCU \(9.204\,\mathrm{s}\)、QUDA \(7.979\,\mathrm{s}\)，比值
  \(0.867\)；这是双 P100 最明显的大格退化项；
\item 同格点 MG-3、trace-off：PyQCU \(2.812\,\mathrm{s}\)、
  QUDA \(4.998\,\mathrm{s}\)，比值 \(1.777\)；
\item V100、\(16\times16\times32\times32\)、c128、MG-2、trace-off：
  PyQCU \(2.487\,\mathrm{s}\)、QUDA \(4.017\,\mathrm{s}\)，比值
  \(1.615\)；
\item 同格点 MG-3：PyQCU \(1.418\,\mathrm{s}\)、QUDA
  \(3.459\,\mathrm{s}\)，比值 \(2.439\)；
\item 全局最大值出现在 V100、\(8^3\times16\)、c128、MG-3、
  trace-off：比值 \(19.587\)，但该结论仅适用于该固定协议。
\end{itemize}

\begin{figure}[htbp]
\centering
\includegraphics[width=0.92\textwidth]{../data/mg_matrix_final_report/mg_speedup_units.pdf}
\caption{按设备与精度统计的 QUDA/PyQCU steady 中位数比值。}
\end{figure}

\subsection{逐层耗时}

trace-on 数据表明，PyQCU 在双 P100 大格的主要成本仍在最粗层
BiCGStab 和 halo/归约延迟。以双 P100 c64、MG-2 trace-on 为例，
PyQCU level-1 coarse solver 约 \(79.66\,\mathrm{s}\)（整次 phase 的
嵌套计时），最细层 pre/post smoother 合计约 \(3.80\,\mathrm{s}\)；
QUDA 的同 phase coarse solver 约 \(28.59\,\mathrm{s}\)。这说明后续
第一优先级仍是粗层流水线，而不是最细层 hopping。

\begin{figure}[htbp]
\centering
\includegraphics[width=0.95\textwidth]{../data/mg_matrix_final_report/mg_stage_breakdown.pdf}
\caption{逐层阶段时间分解（由最终 combined JSON 生成）。}
\end{figure}

\begin{figure}[htbp]
\centering
\includegraphics[width=0.95\textwidth]{../data/mg_matrix_final_report/mg_level_iterations.pdf}
\caption{逐层迭代计数；最细层为 outer FGMRES/GCR，粗层计数独立记录。}
\end{figure}

\subsection{残差}

所有 \(144\) 个 side case 的 true residual 均低于对应 gate。
trace-on 残差序列从 combined JSON 直接提取，不把 Arnoldi 估计值
当作 true residual。图 \ref{fig:residual} 展示各设备、精度、
lattice、MG-2/3 的 trace-on 残差曲线。

\begin{figure}[htbp]
\centering
\includegraphics[width=0.98\textwidth,height=0.78\textheight,keepaspectratio]{../data/mg_matrix_final_report/mg_residual_curves.pdf}
\caption{最细层 outer residual 的迭代曲线；MG-2/3、trace-on。}
\label{fig:residual}
\end{figure}

\begin{figure}[htbp]
\centering
\includegraphics[width=0.92\textwidth]{../data/mg_matrix_final_report/mg_finest_residual.pdf}
\caption{各单元最细层最终真残差与 gate 的对照。}
\end{figure}

\section{正确性与工程验证}

\begin{table}[htbp]
\centering
\small
\caption{最终验证门禁}
\begin{tabularx}{\textwidth}{@{}L{0.30\textwidth}Y@{}}
\toprule
检查 & 结果 \\
\midrule
\code{validate_document} & \(72/72\) 个 combined 文档通过 \\
两侧 config hash & \(72/72\) 完全一致 \\
input bundle hash & \(72/72\) 完全一致 \\
MG level timing conservation & 全部通过，允许的 trace 归约钳零未掩盖失败 \\
分布式 HDF5 round-trip & gauge/source/null \(max\_abs=0\) \\
CUDA/MPI 长期测试 & P100 与 V100 全矩阵均完成，QUDA ghost mapped-memory 失败以重试记录，不改算法协议 \\
测试套件 & \pathref{examples/qcu/dev87/test_bench_strict_protocol.py} 与
  \pathref{examples/qcu/dev87/test_bench_mg_matrix_full.py} 全部通过（最终回归 89 passed） \\
\bottomrule
\end{tabularx}
\end{table}

\section{遗留问题与适用边界}

\begin{enumerate}
\item 双 P100 大格 c64 MG-2 仍为 \(0.867\times\)；粗层
  BiCGStab 的 host-visible reduction 与 halo 延迟没有完全消失；
\item V100 的优势包含设备代际差、单 rank 路径和不同 nu 设置，
  不能外推为任意 GPU/GPU 图上的普遍结论；
\item BiCGstab reference 在全部 24 个 L1 对照中慢于 QUDA，若用户
  关注无 MG 路径，应单独优化该 solver；
\item trace-on 的 stage 时间包含 trace 同步开销，性能比较以
  trace-off 的 steady median 为准；
\item QUDA 在部分 large c128 运行中出现 \code{cudaHostRegister}
  56\,MB mapped-memory 瞬时失败；同一配置重试后成功，记录为
  环境级 ghost-buffer 注册问题，未修改 QUDA 数值算法；
\item 分布式读取已使用 rank-sliced h5py MPI-IO；单 rank 路径仍是
  完整读取，这是为了保持单卡语义和既有 cache 格式；
\item 66/72 等单元在不同 lattice/level 上存在显著方差；报告以
  逐单元 JSON 和 CSV 为准，不用单一平均值替代全部结果。
\end{enumerate}

\section{复现与产物}

\begin{lstlisting}[basicstyle=\ttfamily\small,breaklines=true]
# Build final sm60+sm70 library
QCU_CUDA_ARCHITECTURES='70-real;70-virtual;60-real;60-virtual' \
  bash ./build.sh

# P100 final PyQCU side (example: c64 large)
source examples/qcu/dev87/p100_env.sh
python -B examples/qcu/dev87/bench_mg_matrix_full.py \
  --output-dir data/mg_matrix_final_p100 \
  --cache-root data/mg_matrix_full_20260921_aligned/cache \
  --only-side pyqcu --only-device p100 --execute

# V100 final PyQCU side (c128 uses nu=2; c64 uses nu=1)
CUDA_VISIBLE_DEVICES=0 python -B \
  examples/qcu/dev87/bench_mg_matrix_full.py \
  --output-dir data/mg_matrix_final_v100 \
  --cache-root data/mg_matrix_full_20260921_aligned/cache \
  --only-side pyqcu --only-device v100 --extra-args '--nu 2' \
  --execute

# Full report tables/figures
python -B examples/qcu/dev87/build_mg_report.py \
  --input data/mg_matrix_final_combined/units/*.json \
  --outdir data/mg_matrix_final_report \
  --tex docs/report_multigrid_final_20260921.tex
python -B examples/qcu/dev87/plot_mg_residual_curves.py
\end{lstlisting}

关键产物：
\begin{itemize}
\item 最终 combined：\pathref{data/mg_matrix_final_combined/units/}；
\item 性能统计：\pathref{data/mg_matrix_final_analysis_full.json}；
\item 独立校验：\pathref{data/mg_matrix_final_validation_full.json}；
\item 逐单元 CSV：\pathref{data/mg_matrix_final_report/mg_matrix.csv}；
\item 逐层 CSV：\pathref{data/mg_matrix_final_report/mg_levels.csv}；
\item 图：\pathref{data/mg_matrix_final_report/mg_speedup_units.pdf}、
  \pathref{data/mg_matrix_final_report/mg_stage_breakdown.pdf}、
  \pathref{data/mg_matrix_final_report/mg_level_iterations.pdf}、
  \pathref{data/mg_matrix_final_report/mg_finest_residual.pdf}、
  \pathref{data/mg_matrix_final_report/mg_residual_curves.pdf}。
\end{itemize}

\section{结论}

本轮已经完成非法内存访问 debug、真正的多 rank 输入读取、halo
slot/cache 优化，以及 144 side case 的完整测试。最终事实可以概括为：
\begin{itemize}
\item 正确性：\(72/72\) 个对照全部 fair，144 side case 真残差门禁全过；
\item V100：MG-2/3 \(23/24\) 胜，中位数 \(2.492\times\)；
\item 双 P100：MG-2/3 \(13/24\) 胜，中位数 \(1.063\times\)；大格 MG-2
  仍需继续优化；
\item BiCGstab reference：\(0/24\) 胜，明确保留为未解决短板；
\item 两个“降同步”粗层实验在 large c64 数值失稳后已回退，最终
  版本不包含未经稳定性验证的性能路径。
\end{itemize}

\end{document}
